% =============================================================================
% Template LaTeX Minimalista (Azul & Cinza/Ardósia) — Livro do Minicurso
% II DataPOL · UFPE · 2026
% =============================================================================

% --- Geometria da página -----------------------------------------------------
\usepackage[letterpaper,left=2.6cm,right=2.6cm,top=2.7cm,bottom=3.1cm,headheight=22pt,headsep=0.8cm,footskip=1.25cm]{geometry}
\usepackage{emptypage}

% --- Tipografia (XeLaTeX) ----------------------------------------------------
\usepackage{fontspec}
\defaultfontfeatures{Ligatures=TeX}
\setmainfont{Noto Serif}[Scale=0.96]
\setsansfont{Lato}[
  UprightFont = Lato Regular,
  BoldFont = Lato Bold,
  ItalicFont = Lato Italic,
  BoldItalicFont = Lato Bold Italic
]
\setmonofont{DejaVu Sans Mono}[Scale=0.82]

% Seta Unicode (→) renderizada com glifo matemático nítido
\usepackage{newunicodechar}
\newunicodechar{→}{\ensuremath{\rightarrow}}

\usepackage{setspace}
\setstretch{1.14}
\setlength{\parskip}{0.45em}
\setlength{\parindent}{0pt}

% --- Paleta de cores minimalista (Azul & Cinza) ------------------------------
\usepackage{xcolor}
\usepackage{colortbl}
\definecolor{azulNoite}{HTML}{0F2440}
\definecolor{azulEscuro}{HTML}{1C3F6E}
\definecolor{azulMedio}{HTML}{2B5B9E}
\definecolor{azul}{HTML}{3D7CCC}
\definecolor{azulClaro}{HTML}{7BA4DE}
\definecolor{azulBorda}{HTML}{C9DCF5}
\definecolor{azulFundo}{HTML}{F2F7FD}

\definecolor{tinta}{HTML}{182026}
\definecolor{cinzaEscuro}{HTML}{334155}
\definecolor{cinzaMedio}{HTML}{64748B}
\definecolor{cinzaClaro}{HTML}{94A3B8}
\definecolor{cinzaBorda}{HTML}{CBD5E1}
\definecolor{cinzaFundo}{HTML}{F4F6F9}

\colorlet{corTituloCaixa}{azulEscuro}
\AtBeginDocument{\color{tinta}}

% --- Links e referências cruzadas --------------------------------------------
\AtBeginDocument{%
  \hypersetup{
    colorlinks=true,
    linkcolor=azulEscuro,
    citecolor=azulEscuro,
    urlcolor=azulMedio,
    filecolor=azulMedio
  }%
}

% --- Folha de rosto interna minimalista (\maketitle) -------------------------
\makeatletter
\renewcommand{\maketitle}{%
  \begin{titlepage}%
    \sffamily
    \vspace*{2.5cm}%
    {\footnotesize\bfseries\color{azul}\MakeUppercase{II DataPOL · UFPE · 2026}\par}%
    \vspace{0.8em}%
    {\Huge\bfseries\color{azulNoite}\@title\par}%
    \vspace{0.6em}%
    {\color{azul}\rule{2.4cm}{2.5pt}}{\color{cinzaBorda}\rule{5.6cm}{0.6pt}}\par
    \vspace{0.8em}%
    \@ifundefined{@subtitle}{}{{\Large\color{cinzaEscuro}\@subtitle\par}}%
    \vfill
    {\color{cinzaBorda}\hrule height 0.5pt}\vspace{1em}%
    {\def\and{\par\vspace{0.25em}}\normalsize\color{tinta}\@author\par}%
    \vspace{0.6em}%
    {\small\color{cinzaMedio}\@date\par}%
    \vspace*{1cm}%
  \end{titlepage}%
}
\makeatother

% --- Cabeçalhos de capítulos e seções (titlesec) -----------------------------
\usepackage{titlesec}

\titleformat{\chapter}[display]
  {\normalfont\sffamily}
  {\footnotesize\bfseries\color{azul}\MakeUppercase{\chaptertitlename}\ \thechapter}
  {0.35em}
  {\huge\bfseries\color{azulNoite}}
  [\vspace{0.45em}{\color{azul}\rule{2.2cm}{2.2pt}}{\color{cinzaBorda}\rule{\dimexpr\textwidth-2.2cm\relax}{0.6pt}}]

\titleformat{name=\chapter,numberless}[display]
  {\normalfont\sffamily}
  {}
  {0pt}
  {\huge\bfseries\color{azulNoite}}
  [\vspace{0.45em}{\color{azul}\rule{2.2cm}{2.2pt}}{\color{cinzaBorda}\rule{\dimexpr\textwidth-2.2cm\relax}{0.6pt}}]

\titlespacing*{\chapter}{0pt}{-10pt}{1.8em}

\titleformat{\section}
  {\normalfont\sffamily\Large\bfseries\color{azulEscuro}}
  {\textcolor{azul}{\thesection}}
  {0.6em}
  {}
  [\vspace{0.2em}{\color{cinzaBorda}\hrule height 0.4pt}]
\titlespacing*{\section}{0pt}{1.6em}{0.6em}

\titleformat{\subsection}
  {\normalfont\sffamily\large\bfseries\color{cinzaEscuro}}
  {\textcolor{cinzaMedio}{\thesubsection}}
  {0.55em}
  {}
\titlespacing*{\subsection}{0pt}{1.2em}{0.4em}

\titleformat{\subsubsection}
  {\normalfont\sffamily\normalsize\bfseries\color{cinzaEscuro}}
  {\thesubsubsection}
  {0.5em}
  {}

% --- Cabeçalho e rodapé das páginas (fancyhdr) -------------------------------
\usepackage{fancyhdr}
\pagestyle{fancy}
\fancyhf{}
\fancyhead[LE]{\sffamily\footnotesize\color{cinzaMedio}\nouppercase{\leftmark}}
\fancyhead[RO]{\sffamily\footnotesize\color{cinzaMedio}\nouppercase{\rightmark}}
\fancyfoot[LE,RO]{\sffamily\footnotesize\bfseries\color{azulEscuro}\thepage}
\fancyfoot[LO,RE]{\sffamily\scriptsize\color{cinzaClaro}II DataPOL · IA como Instrumento de Medição}
\renewcommand{\headrulewidth}{0.4pt}
\renewcommand{\headrule}{\hbox to\headwidth{\color{cinzaBorda}\leaders\hrule height \headrulewidth\hfill}}
\renewcommand{\footrulewidth}{0pt}

\fancypagestyle{plain}{%
  \fancyhf{}%
  \fancyfoot[LE,RO]{\sffamily\footnotesize\bfseries\color{azulEscuro}\thepage}%
  \renewcommand{\headrulewidth}{0pt}%
}

% --- Caixas de destaque (tcolorbox + filtro caixa.lua) -----------------------
\usepackage[most]{tcolorbox}

\newtcolorbox{caixanota}{
  enhanced,
  breakable,
  lines before break=4,
  pad at break*=2.5mm,
  colback=azulFundo,
  colframe=azulBorda,
  borderline west={3.5pt}{0pt}{azulEscuro},
  boxrule=0.5pt,
  arc=3pt,
  left=11pt, right=11pt, top=9pt, bottom=9pt,
  before skip=1.1em, after skip=1.1em,
  before upper={\colorlet{corTituloCaixa}{azulEscuro}\setlength{\parskip}{0.4em}}
}

\newtcolorbox{caixaaviso}{
  enhanced,
  breakable,
  lines before break=4,
  pad at break*=2.5mm,
  colback=cinzaFundo,
  colframe=cinzaBorda,
  borderline west={3.5pt}{0pt}{cinzaEscuro},
  boxrule=0.5pt,
  arc=3pt,
  left=11pt, right=11pt, top=9pt, bottom=9pt,
  before skip=1.1em, after skip=1.1em,
  before upper={\colorlet{corTituloCaixa}{cinzaEscuro}\setlength{\parskip}{0.4em}}
}

\newtcolorbox{caixaia}{
  enhanced,
  breakable,
  lines before break=4,
  pad at break*=2.5mm,
  colback=azulFundo,
  colframe=azulBorda,
  borderline west={3.5pt}{0pt}{azul},
  boxrule=0.5pt,
  arc=3pt,
  left=11pt, right=11pt, top=9pt, bottom=9pt,
  before skip=1.1em, after skip=1.1em,
  before upper={\colorlet{corTituloCaixa}{azulMedio}\setlength{\parskip}{0.4em}}
}

% Citações (blockquote) em estilo minimalista
\AtBeginDocument{%
  \renewenvironment{quote}{%
    \begin{tcolorbox}[
      enhanced, breakable,
      lines before break=3,
      colback=cinzaFundo, colframe=cinzaFundo,
      borderline west={2.8pt}{0pt}{azulClaro},
      boxrule=0pt, sharp corners,
      left=10pt, right=10pt, top=6pt, bottom=6pt,
      before skip=0.9em, after skip=0.9em,
      fontupper=\small\color{cinzaEscuro}
    ]%
  }{%
    \end{tcolorbox}%
  }%
}

% --- Blocos de código e realce em tons de Azul & Cinza -----------------------
\usepackage{fvextra}
\makeatletter
\AtBeginDocument{%
  \RecustomVerbatimEnvironment{Highlighting}{Verbatim}{commandchars=\\\{\},fontsize=\footnotesize,breaklines,breakanywhere}%
  \let\verbatim\relax\let\endverbatim\relax
  \DefineVerbatimEnvironment{verbatim}{Verbatim}{fontsize=\footnotesize,breaklines,breakanywhere,frame=leftline,framerule=2pt,rulecolor=\color{cinzaBorda},framesep=7pt}%
  \@ifundefined{Shaded}{}{%
    \renewenvironment{Shaded}{%
      \begin{tcolorbox}[
        enhanced, breakable,
        lines before break=3,
        colback=cinzaFundo, colframe=cinzaBorda,
        borderline west={2.5pt}{0pt}{azulEscuro},
        boxrule=0.4pt, arc=2.5pt,
        left=8pt, right=8pt, top=5pt, bottom=5pt,
        before skip=0.8em, after skip=0.8em
      ]%
    }{%
      \end{tcolorbox}%
    }%
  }%
  % Paleta de realce de sintaxe em azul e cinza
  \@ifundefined{KeywordTok}{}{%
    \renewcommand{\KeywordTok}[1]{\textcolor[HTML]{1C3F6E}{\textbf{#1}}}%
    \renewcommand{\ControlFlowTok}[1]{\textcolor[HTML]{1C3F6E}{\textbf{#1}}}%
    \renewcommand{\OperatorTok}[1]{\textcolor[HTML]{1C3F6E}{#1}}%
    \renewcommand{\FunctionTok}[1]{\textcolor[HTML]{2B5B9E}{#1}}%
    \renewcommand{\BuiltInTok}[1]{\textcolor[HTML]{2B5B9E}{#1}}%
    \renewcommand{\StringTok}[1]{\textcolor[HTML]{2B6CB0}{#1}}%
    \renewcommand{\CharTok}[1]{\textcolor[HTML]{2B6CB0}{#1}}%
    \renewcommand{\SpecialStringTok}[1]{\textcolor[HTML]{2B6CB0}{#1}}%
    \renewcommand{\VerbatimStringTok}[1]{\textcolor[HTML]{2B6CB0}{#1}}%
    \renewcommand{\DecValTok}[1]{\textcolor[HTML]{1E4E8C}{#1}}%
    \renewcommand{\BaseNTok}[1]{\textcolor[HTML]{1E4E8C}{#1}}%
    \renewcommand{\FloatTok}[1]{\textcolor[HTML]{1E4E8C}{#1}}%
    \renewcommand{\ConstantTok}[1]{\textcolor[HTML]{1E4E8C}{#1}}%
    \renewcommand{\AttributeTok}[1]{\textcolor[HTML]{2B5B9E}{#1}}%
    \renewcommand{\CommentTok}[1]{\textcolor[HTML]{64748B}{\textit{#1}}}%
    \renewcommand{\DocumentationTok}[1]{\textcolor[HTML]{64748B}{\textit{#1}}}%
    \renewcommand{\AnnotationTok}[1]{\textcolor[HTML]{64748B}{\textit{#1}}}%
    \renewcommand{\CommentVarTok}[1]{\textcolor[HTML]{64748B}{\textit{#1}}}%
    \renewcommand{\NormalTok}[1]{\textcolor[HTML]{182026}{#1}}%
    \renewcommand{\VariableTok}[1]{\textcolor[HTML]{182026}{#1}}%
    \renewcommand{\OtherTok}[1]{\textcolor[HTML]{334155}{#1}}%
  }%
}
\makeatother

% Caminhos de arquivo em \texttt podem quebrar depois de / . e _ (só dentro de \texttt)
\newXeTeXintercharclass\classebarra
\XeTeXcharclass`\/=\classebarra
\XeTeXcharclass`\.=\classebarra
\XeTeXinterchartoks\classebarra 0={\allowbreak}
\let\textttorig\texttt
\renewcommand{\texttt}[1]{\textttorig{\color{azulEscuro}\XeTeXinterchartokenstate=1\def\_{\textunderscore\allowbreak}#1}}

% --- Tabelas e legendas ------------------------------------------------------
\usepackage{caption}
\captionsetup{
  font={small,sf,color=cinzaEscuro},
  labelfont={bf,sf,color=azulEscuro},
  labelsep=period,
  skip=6pt
}
\AtBeginEnvironment{longtable}{\small\sffamily\arrayrulecolor{cinzaBorda}}
\setlength{\LTpre}{0.9em}
\setlength{\LTpost}{0.9em}
